% Escreva aqui a solução da questão 5.


\textcolor{red}{ VERSÃO 1 } %Letícia

%Inicialmente, analisaremos a condição de existência do radicando para a radiciação com índice par (índice igual a 2): $2x + 1  + x^{2} \geq 0$.

%Calculando o $\Delta$: $\Delta = b^{2} - 4ac \implies \Delta = 2^{2} -4 \cdot 1 \cdot 1 \implies \Delta = 0$. Logo, teremos duas raízes reais e iguais. Como $a = 1 > 0$, a concavidade da parábola é voltada para cima. Dessa maneira, a expressão $2x + 1  + x^{2}$ será sempre maior ou igual a zero para todo $x \in \mathbb{R}$.

%Agora, resolvemos a equação:
Segue a solução da equação:
\begin{align*} %Lembrar de não pular linha para não quebrar o alinhamento do align
    \sqrt{2x + 1  + x^{2}} = 1 & \implies (\sqrt{2x + 1  + x^{2}})^{2} = 1^{2} \\
    & \iff 2x + 1 + x^{2} = 1 \\
    & \iff x^{2} + 2x = 0 \\
    & \iff x(x+2) = 0. 
\end{align*}

Então, temos dois casos: $x = 0$ ou $x + 2 = 0 \iff x = -2$.

Verificando a validade das soluções na equação original:

\begin{itemize}
    \item $x = 0$: $\sqrt{2 \cdot 0 + 1 + 0^{2}} = 1 \iff \sqrt{1} = 1$. Então, $x = 0$ é solução.
    \item $x = -2$: $\sqrt{2 \cdot (-2) + 1 + (-2)^{2}} = 1 \iff \sqrt{1} = 1$. Então, $x = -2$ é solução.
\end{itemize}
Logo, o conjunto solução é: $S = \{-2, 0 \}$.


\bigskip


\textcolor{red}{ VERSÃO 2} %Letícia

\begin{align*} %Lembrar de não pular linha para não quebrar o alinhamento do align
    \sqrt{2x + 1  + x^{2}} = 1 & \iff \sqrt{(x+1)^{2}} = 1^{2} \\
    & \iff |x + 1| = 1. 
\end{align*}
Como $$|x + 1| = \left\{
            \begin{array}{r l}
                 x + 1, & \text{ se } x+1 \geq 0 \iff x \geq -1\\
                 -(x +1), & \text{ se } x+1 < 0 \iff x < -1 
            \end{array},
            \right .$$
        então:
    \begin{itemize}
        \item Caso $x \geq -1$: $ |x + 1| = 1 \iff x + 1 = 1 \iff x = 0$. 
        
        Como $0 \geq -1$, então zero é uma solução da equação.
        \item Caso $x < -1$: $ |x + 1| = 1 \iff -(x +1) = 1 \iff x = -2$. 
        
        Como $-2 < -1$, então $-2$ é uma solução da equação.
    \end{itemize}
    Logo, o conjunto solução é: $S = \{-2, 0 \}$.





%\bigskip

%\textcolor{red}{ VERSÃO 3} % autor




%\bigskip





\bigskip

\textcolor{red}{CHAVE DE PONTUAÇÃO:
\begin{itemize}
    \item  Solução da equação SEM os testes da solução \emph{(2,0 pontos)};
    \item  Solução da equação COM os testes da solução \emph{(3,0 ponto)}.
\end{itemize} }